% ---- score2 standard preamble (from: python3 score2_tools.py preamble) ----
\documentclass[10pt]{extarticle}
\usepackage[a4paper,margin=1.5cm]{geometry}
\usepackage[T1]{fontenc}
\usepackage{amsmath,amssymb}
\usepackage{graphicx}
\usepackage[export]{adjustbox}
\usepackage{multicol}
\usepackage{enumitem}
\usepackage{booktabs,array,tabularx}
\setlength{\columnsep}{0.7cm}
\setlength{\parindent}{0pt}
\setlength{\parskip}{3pt}
\pagestyle{empty}
\makeatletter
% figure registry (figdefs.tex is generated by the extract step)
\newcommand{\qfigdef}[4]{\@namedef{qf@file@#1}{#2}\@namedef{qf@w@#1}{#3}\@namedef{qf@wide@#1}{#4}}
\newcommand{\qf@check}[1]{\@ifundefined{qf@file@#1}{\PackageError{score2}{Unknown figure id #1}{}}{}}
\newcommand{\qf@img}[4][]{\includegraphics[width=#2\dimexpr\@nameuse{qf@w@#4}pt\relax,max width=#3,max height=0.45\textheight,#1]{\@nameuse{qf@file@#4}}}
% Full-width questions: finalize wraps every question that contains a wide figure in \qfullbegin...\qfullend,
% so the WHOLE question (stem, figures, options) is set across the page in reading order, then the two columns
% resume. The question's head - its number, stem and first figure (\qfullsplit marks the end of the head) -
% is kept on one page, so a number is never left at a page foot with its figure overleaf; the rest may break
% across pages normally, so no big gaps are left. Without a \qfullsplit, the whole question is kept together
% if it is under 45% of a page. Never write these yourself.
\newif\ifqf@full
\newif\ifqf@boxed
\newsavebox{\qf@box}
\newcommand{\qf@open}{\global\qf@fulltrue\par\global\qf@boxedtrue
  \setbox\qf@box\vbox\bgroup\hsize\textwidth\linewidth\textwidth}
\newcommand{\qf@emit}[1]{\par\egroup\global\qf@boxedfalse
  \ifdim\dimexpr\ht\qf@box+\dp\qf@box\relax<#1\textheight\noindent\usebox{\qf@box}\par\else\unvbox\qf@box\fi}
\newcommand{\qfullsplit}{\ifqf@boxed\qf@emit{0.85}\fi}
\newcommand{\qf@close}{\ifqf@boxed\qf@emit{0.45}\fi\global\qf@fullfalse}
\newcommand{\qfullbegin}{\end{multicols}\qf@open}
\newcommand{\qfullend}{\qf@close\begin{multicols}{2}}
\newcommand{\qfullstart}{\qf@open}
\newcommand{\qfulldone}{\qf@close}
% \qfig[scale]{id}: block figure, centred, directly where it is written (size set by the tool).
% Fallback if a wide figure is not inside a full-width question: break out of the columns in place - no floats.
\newcommand{\qfig}[2][1]{\qf@check{#2}%
  \ifqf@full\qf@blk{#1}{#2}\else
  \ifnum\@nameuse{qf@wide@#2}=1 \expandafter\@firstoftwo\else\expandafter\@secondoftwo\fi
  {\end{multicols}\begin{center}\qf@img{#1}{\textwidth}{#2}\end{center}\begin{multicols}{2}}%
  {\qf@blk{#1}{#2}}\fi}
\newcommand{\qf@blk}[2]{\par\smallskip{\centering\qf@img{#1}{\linewidth}{#2}\par}\smallskip}
% \qfigi[scale]{id}: inline figure (inside option lists / tables / minipages); never breaks columns.
\newcommand{\qfigi}[2][1]{\qf@check{#2}\qf@img[valign=c]{#1}{\linewidth}{#2}}
\makeatother
% structure macros - lint reads these to check completeness and figure placement
\newcommand{\papersec}[1]{\par\bigskip{\centering\Large\bfseries PAPER-#1\par}\medskip}
\newcommand{\partsec}[2]{\par\medskip\noindent{\large\bfseries PART-#1: #2}\par\smallskip}
\newcommand{\sect}[1]{\par\smallskip\noindent\textbf{#1}\par}
\newcommand{\qn}[1]{\par\smallskip\noindent\textbf{#1.}~}
\newcommand{\paraq}[2]{\par\smallskip\noindent\textbf{Paragraph for Questions #1 and #2}\par}
\newenvironment{qoptions}{\begin{itemize}[leftmargin=1.8em,itemsep=1pt,topsep=2pt,label={}]}{\end{itemize}}
\newcommand{\opt}[2]{\item[(#1)] #2}
% figure registry (generated by score2_tools.py finalize)
\qfigdef{p1_phy_s1_q06_1}{Score-I - Test 26 (18 Dec 2022)-img/p1_phy_s1_q06_1.png}{39.0}{0}
\qfigdef{p1_phy_s3_para3_4_1}{Score-I - Test 26 (18 Dec 2022)-img/p1_phy_s3_para3_4_1.png}{164.6}{0}
\qfigdef{p1_phy_s4_q07_1}{Score-I - Test 26 (18 Dec 2022)-img/p1_phy_s4_q07_1.png}{108.7}{0}
\qfigdef{p1_chem_s1_q03_1}{Score-I - Test 26 (18 Dec 2022)-img/p1_chem_s1_q03_1.png}{117.1}{0}
\qfigdef{p1_chem_s3_para5_6_1}{Score-I - Test 26 (18 Dec 2022)-img/p1_chem_s3_para5_6_1.png}{60.8}{0}
% ---- body skeleton ----
% \begin{document}
% \begin{center}{\Large\bfseries JEE (Main+Advanced) Score-II --- Test NN (DD Mon YYYY)}\end{center}
% \begin{multicols}{2}
%   \papersec{1}   (only for two-paper tests)
%   \partsec{1}{PHYSICS} \sect{Section-I (i): Single Correct} \qn{1} ... \qfig{p1_phy_q01_1} ...
% \end{multicols}
% \end{document}

\begin{document}
\begin{center}{\Large\bfseries JEE (Main+Advanced) Score-I --- Test 26 (18 Dec 2022)}\end{center}
\raggedcolumns
\setlength{\tabcolsep}{3pt}
\begin{multicols}{2}
\papersec{1}
\partsec{1}{PHYSICS}
\sect{SECTION-I(i)}
\qn{1} A certain microwave oven projects 1.77 kW of microwaves onto a 100 cm by 20 cm area.
\begin{qoptions}
\opt{A}{The intensity of microwaves is $8.85\times10^{3}$ W/m$^2$}
\opt{B}{The peak intensity of microwaves is $8.85\times10^{3}$ W/m$^2$}
\opt{C}{The square of peak electric field strength in the waves is $6.66\times10^{6}$ N$^2$/c$^2$}
\opt{D}{The peak value of magnetic energy density $2.95\times10^{-5}$ J/m$^3$}
\end{qoptions}

\qn{2} Mark the correct statements about polarisation of light.
\begin{qoptions}
\opt{A}{If a polaroid is kept in front of you and rotated while looking at the light reflected from water or glass, you will notice that light gets bright and completely dark sequencially.}
\opt{B}{At Brewster's angle of incidence the plane of vibration of electric field vector of reflected light is perpendicular to plane of reflection}
\opt{C}{A polaroid consists of long chain molecules alinged in direction parallel to the axis of polaroid which absorb the electromagnetic energy of light.}
\opt{D}{If $\theta_1$ is the Brewster's angle for light reflected from the top of an interface between two substances and $\theta_2$ is the Brewster's angle for light reflected from below, then $\theta_1+\theta_2=90^\circ$.}
\end{qoptions}

\qn{3} In a single-slit diffraction pattern the central fringe is 450 times as wide as the slit. The screen is 18000 times farther from the slit than the slit is wide. The wave length of light is $\lambda$ and width of slit is `W'. [$\pi^2=10$]
\begin{qoptions}
\opt{A}{The ratio $\dfrac{W}{\lambda}$ is 80.}
\opt{B}{The distance of second minima from the point of maximum intensity is 450 W}
\opt{C}{The distance of second maxima from the point of maximum intensity is 450 W}
\opt{D}{If the angular position of the point where intensity is 0.9 times the maximum intensity is $\theta$, then $\sin\theta\approx\dfrac{1}{480}$ .}
\end{qoptions}

\qn{4} Mark the correct statements of the following :
\begin{qoptions}
\opt{A}{According to Rayleigh's criterion for resolution, two images would be just resolved when central maxima of one image coincides with central maxima of the other.}
\opt{B}{The resolving power of a microscope is inversely proportional to the focal length of the objective lens of microscope.}
\opt{C}{The resolving power of a telescope is directly proportional to the aperture diameter of the objective lens.}
\opt{D}{For an eagle flying at an altitude of 100 m having pupil diameter of 6 mm, two point objects at a separation of 20 mm can be resolved considering wave length of light $\lambda=600$ nm.}
\end{qoptions}

\qn{5} A source of sound radiates uniformly in all directions. Along a radial line from the source locate two points separated by 3 m such that intensity at the nearer point is 4dB above that of more distance point. [given $10^{0.2}=1.6$]
\begin{qoptions}
\opt{A}{Ratio of intensities at the nearer point to the distance point is 1.6}
\opt{B}{Ratio of distance of distant point to the nearer point from the source is 1.6}
\opt{C}{Distance of the nearer point from source is 5 m.}
\opt{D}{Distance of the distant point from source is 8 m.}
\end{qoptions}

\qn{6} A composite wire is made by joining 3 wires of length $\ell$ and linear mass density $\mu$, $2\mu$ and $3\mu$ respectively as shown in figure. Wire AB has the least density followed by BC and CD.

A quick transverse pulse is given at the end D at $t=0$.
\qfig{p1_phy_s1_q06_1}
\begin{qoptions}
\opt{A}{The ratio of time taken by pulse to travel the DC, DB and DA is $1:\sqrt2:\sqrt3$}
\opt{B}{The ratio of time taken by pulse to travel the DC, CB and BA is $1:\dfrac{1}{\sqrt2}:\dfrac{1}{\sqrt3}$}
\opt{C}{The ratio of length of the pulse just above B and just below C is $\sqrt5:1$}
\opt{D}{Acceleration of the pulse travelling upwards is in the ratio of 3 : 5 : 6 in the strings CD, BC and AB respectively.}
\end{qoptions}

\sect{SECTION-I(ii)}
\paraq{7}{8}
A detector of radio waves in a radio astronomical observatory is placed on the sea beach at a height $h=2.10$ m above the sea level. After the rise of a star, radiating electromagnetic waves of wavelength $\lambda=21$ cm, above the horizon the detector registers series of alternating maxima and minima. The registered signal is proportional to the intensity of detected waves. The detector registers waves with electric vector vibrating parallel to the sea surface.

\qn{7} If the angle between the star and horizon in the moment the detector resistors maxima for the second time is $\theta_0$, the value of $\sin\theta_0$ is given by $N\times10^{-3}$. Find the value of $N$.
\begin{qoptions}
\opt{A}{55}
\opt{B}{65}
\opt{C}{75}
\opt{D}{85}
\end{qoptions}

\qn{8} At reflection of the electromagnetic wave on water surface the ratio of the intensities of the electric field of the reflected and ($E_r$) and incident ($E_i$) wave follows the law :
\[\frac{E_r}{E_i}=\frac{n-\cos i}{n+\cos i}\]
where $n$ is the refractive - index and $i$ is the incident angle of wave. For the surface ``air-water'' for $\lambda=21$ cm, the refractive index $n=9$

If the signal ratio of the first maximum to the next minimum is $R\times10^{4}$, find the value of $R$.
\begin{qoptions}
\opt{A}{3.25}
\opt{B}{4.25}
\opt{C}{5.25}
\opt{D}{None of these}
\end{qoptions}

\paraq{9}{10}
Three successive resonance frequencies for a certain string are 175, 245 and 315 Hz. The length of the string is 12.87 m. The separation between the nodes in the respective modes are $\alpha$, $\beta$ and $\gamma$ in meters. the fundamental frequency of the this string is $f_0$ mass per unit length of the string is $10^{-3}$ kg/m.

\qn{9} Find the sum of $\alpha$, $\beta$ and $\gamma$.
\begin{qoptions}
\opt{A}{6.56}
\opt{B}{7.18}
\opt{C}{6.18}
\opt{D}{11.68}
\end{qoptions}

\qn{10} The same string resonates with the frequency $f_0$, but has two nodes and two antinodes on the string. Find the tension in the string in newtons.
\begin{qoptions}
\opt{A}{360.72}
\opt{B}{260.72}
\opt{C}{266.72}
\opt{D}{366.72}
\end{qoptions}

\sect{SECTION-II(i)}
\paraq{1}{2}
A massless speaker emitting sound at 600 Hz is placed on top of a mass connected to a spring. The mass spring system is set in an oscillatory state, vibrating horizontally at 4 Hz with amplitude $A$. The velocity of sound is given to be 330 m/s.

\qn{1} If the difference between highest and lowest frequencies heard by a stationary observer on the line of motion is 2 Hz, find the amplitude $A$ in cm.

\qn{2} If the above spring and block are cut in two equal halves speaker is set on one of the halves connected to one half spring and set in oscillation with same amplitude, find the maximum frequency of sound heard by observer rounded to closest integer value.

\paraq{3}{4}
The power radiated by sun is $4.0\times10^{26}$ W. The earth orbits the sun in a nearly circular orbit of radius $1.5\times10^{11}$ m. The earth's axis is tilted by $27^\circ$ relative to the plane of the orbit, so sunlight does not strike perpendicularly. $\left[\cos27^\circ=0.9,\ \pi=\dfrac{22}{7}\right]$
\qfig{p1_phy_s3_para3_4_1}

\qn{3} Find the power (in W) that strikes a 0.75 m$^2$ patch of flat land at the equator at point Q ?

\qn{4} Find the average value of total energy density `u' of the radiation incident at point Q on earth in $\mu$Jm$^{-3}$.

\paraq{5}{6}
A convex lens of focal length 20 cm is placed after a slit of width 0.6 mm. A plane wave of wavelength 6000 \AA\ falls normally on the slit. A screen is placed at the focal plane of the convex lens.

\qn{5} Find the separation between the second minima on either side of the central maximum (in mm)

\qn{6} Find the ratio of the intensity of the principal maximum to the first maximum on either side of principal maximum. Consider secondary maximas to be placed right at the middle of minimas. [Take $\pi^2=10$]

\sect{SECTION-II(ii)}
\qn{7} In a young's double slit experiment a parallel beam of mono chromatic light of wavelength $\lambda=600$ nm is incident on the plane of slits making an angle $\theta$ with the normal to the plane ($\sin\theta=0.03$). The separation between the slits is $d=0.1$ mm. When a thin film is introduced in front of the the lower slit, three minimas and two maximas go past the center of the screen `O' and the intensity at O reduces to one-fourth of the intensity before. If the refractive index of the material of thin film is 1.2, find the width of the film in $\mu$m.
\qfig{p1_phy_s4_q07_1}

\qn{8} Unpolarised light of intensity $I_0$ is incident on a polariser which is followed by a series of polarisers. Each polariser has its transmission axis rotated by $45^\circ$ relative to the transmission axis of the previous polariser. What is the minimum number of polarisers needed for the light reaching the photo cell to have an intensity that is less than $\dfrac{I_0}{100}$.

\qn{9} A musical instrument has a spectrum of frequencies ranging from 100 Hz to 5000 Hz. The ear canal is a tube of length 3 cm and is open at one end and closed at the other. Find the frequency that will be best heard by ears in Hz. (take speed of sound = 330 m/s)
\partsec{2}{CHEMISTRY}
\sect{SECTION-I(i)}
\qn{1} Which point defect in crystals may reduce the crystal density ?
\begin{qoptions}
\opt{A}{Schottky defect}
\opt{B}{Frenkel defect}
\opt{C}{Metal excess defect}
\opt{D}{Metal deficiency defect}
\end{qoptions}

\qn{2} Under conditions of same temperature and pressure, mark correct relationship.
\begin{qoptions}
\opt{A}{Density of water vapour is more than N$_2$.}
\opt{B}{Density of H$_2$ is more than water vapour.}
\opt{C}{CO$_2$ is denser than air.}
\opt{D}{Order of density : H$_2<$ D$_2<$ T$_2$}
\end{qoptions}

\qn{3} A generic phase diagram of a pure substance is drawn :
\qfig{p1_chem_s1_q03_1}
Choose correct statement :
\begin{qoptions}
\opt{A}{At C, surface tension of liquid becomes zero.}
\opt{B}{At X, pure substance exists in solid, liquid and gaseous state simultaneously}
\opt{C}{XB represents variation of boiling point of the substance, with pressure.}
\opt{D}{At low T and high P, substance exhibits ideal gas behaviour}
\end{qoptions}

\qn{4} Sodium hydroxide cannot be used as a primary standard for acid base titration because
\begin{qoptions}
\opt{A}{it is corrosive in nature and reacts with glass.}
\opt{B}{the dissolution of sodium hydroxide in water is highly exothermic, hence its concentration changes on dissolution.}
\opt{C}{it is hygroscopic and also reacts with atmospheric carbon dioxide.}
\opt{D}{Hydroxides cannot be used as a primary standard.}
\end{qoptions}

\qn{5} The following equilibrium exists in saturated solution of NH$_4$Cl

$\mathrm{NH_4Cl_{(s)}}\rightleftharpoons\mathrm{NH^+_{4(aq)}}+\mathrm{Cl^-_{(aq)}}$

$\Delta H_{25^\circ C}=3.5$ kcal mol$^{-1}$

A change that will not shift the equilibrium to the right is
\begin{qoptions}
\opt{A}{decrease in temperature}
\opt{B}{increase in temperature}
\opt{C}{addition of NH$_4$Cl crystals to the reaction mixture.}
\opt{D}{addition of NH$_4$OH solution to the reaction mixture.}
\end{qoptions}

\qn{6} Wavelength associated with electron in motion
\begin{qoptions}
\opt{A}{increases with decrease in the speed of electrons}
\opt{B}{remains same irrespective of speed of electron}
\opt{C}{decreases with increase in the speed of electron}
\opt{D}{changes with the atomic number of atom to which it belongs}
\end{qoptions}

\sect{SECTION-I(ii)}
\paraq{7}{8}
The wave function for an atomic orbital of single electron atom or ion is
\[\begin{aligned}\Psi(r,\theta,\phi)={}&\frac{2}{3}\left(\frac{Z}{3a_0}\right)^{3/2}(1-\sigma)\\&\times(12-8\sigma+\sigma^2)\,\sigma\, e^{-\sigma/2}\cos\theta.\end{aligned}\]
Where $\sigma=\dfrac{2Zr}{na_0}$, $a_0=0.529$ \AA.

All other parameters have their usual meaning.

\qn{7} $(n+\ell)$ Value of the given atomic orbital will be
\begin{qoptions}
\opt{A}{2}
\opt{B}{3}
\opt{C}{4}
\opt{D}{6}
\end{qoptions}

\qn{8} Find the value of $m_\ell$ for the given orbital (Assuming $z$-axis is the internuclear axis)
\begin{qoptions}
\opt{A}{0}
\opt{B}{2}
\opt{C}{1}
\opt{D}{$-2$}
\end{qoptions}

\paraq{9}{10}
The reaction : $\mathrm{A(g)}\longrightarrow 2\mathrm{B(g)}+\mathrm{C}(\ell)$ obeys first order kinetics. 0.50 moles of A(g) is taken in an empty vessel of 16.4 L capacity and maintained at $27^\circ$C, at which A(g) start decomposing. In first 30 min, moles of B(g) becomes 0.75.

[Given : Vapour pressure of C($\ell$) at $27^\circ$C $=0.6$ atm, $R=0.082$ L-atm / K-mol]

\qn{9} Total pressure of system at $t=30$ min. is :
\begin{qoptions}
\opt{A}{1.875 atm}
\opt{B}{1.3125 atm}
\opt{C}{0.1875 atm}
\opt{D}{1.0312 atm}
\end{qoptions}

\qn{10} The time at which first drop of liquid `C' will form (log 2 = 0.3)
\begin{qoptions}
\opt{A}{30 min}
\opt{B}{35 min}
\opt{C}{15 min}
\opt{D}{60 min}
\end{qoptions}

\sect{SECTION-II(i)}
\paraq{1}{2}
Following equilibria are established on mixing two gases A$_2$ and C.

(i) $3\mathrm{A_{2(g)}}\rightleftharpoons\mathrm{A_{6(g)}}$, \ $K_p=1.6$ atm$^{-2}$

(ii) $\mathrm{A_{2(g)}}+\mathrm{C_{(g)}}\rightleftharpoons\mathrm{A_2C_{(g)}}$

Given that the total pressure is 1.4 atm and partial pressure of A$_6$ is 0.2 atm,

If A$_2$ and C are mixed in 2 : 1 molar ratio, then

\qn{1} Calculate equilibrium partial pressure of C.

\qn{2} The value of $K_p$ for the reaction (ii) (in atm$^{-1}$) is.

\paraq{3}{4}
CO$_2$ gas (1.00 mole) at 400 K occupies 600 ml.

[For CO$_2$ ; $a=3.60$ L$^2$atm mol$^{-2}$; $b=0.04$ L mol$^{-1}$; $R=0.08$ L atm k$^{-1}$ mol$^{-1}$].

\qn{3} Calculate the pressure (in atm) exerted by the gas using ideal gas equation.

\qn{4} Calculate the pressure (in atm) exerted by the gas using Vander waal's gas equation.

\paraq{5}{6}
The given figure can be a plane projection in the cubic crystal packing of identical particles in a solids. The given figure may or may not be a square. Consider atoms to be perfectly circular/spherical as the case may be.
$\left(\text{Packing fraction}=\dfrac{\text{Packing efficiency}}{100}\right)$
\qfig{p1_chem_s3_para5_6_1}

\qn{5} Calculate packing efficiency, if the packing is in 2D.

\qn{6} Calculate maximum packing fraction, if the packing is in 3D. (Express your answer as nearest integer)

\sect{SECTION-II(ii)}
\qn{7} In the reduction of nitric oxide, 50\% of reaction was completed in 100 s, when initial pressure was 360 mm Hg and in 144 s, when initial pressure was 300 mm Hg. Find order of the reaction.

\qn{8} What is the pH at $25^\circ$C of the solution obtained by dissolving a 5.00 g tablet of aspirin in 0.500 L of water ?

The 5 g tablet contains 0.30 g of aspirin (HC$_9$H$_7$O$_4$). The acid is monoprotic and $K_a=\left(\dfrac{10^{-3}}{3}\right)$ at $25^\circ$C.

Write nearest integral value of pH.\\$\left(\sqrt{41}=6.4,\ \sqrt{123}=11.1\right)$

(log 3 = 0.48, log 2 = 0.30, log 7 = 0.85)

\qn{9} If $pd$ versus $p$ (where $p$ in atm and $d$ in g/L) is plotted for He gas at a particular temperature. If $\left[\dfrac{d}{dp}(pd)\right]_{p=8.21\ \text{atm}}=5$, what is the temperature in Kelvin ? Assume ideal gas behaviour. ($R=0.0821$ L-atm/K-mol)
\partsec{3}{MATHEMATICS}
\sect{SECTION-I(i)}
\qn{1} If value of definite integral\[\int_0^{2\pi}\max\left(2\sin x+\cos2x,\ \dfrac{8x-15\pi}{\pi}\right)dx\]is `A' then -

(where $[.]\to$ represent G.I.F.)
\begin{qoptions}
\opt{A}{$A>0$}
\opt{B}{$A<0$}
\opt{C}{$[A]=0$}
\opt{D}{$[A]=-2$}
\end{qoptions}

\qn{2} Let us consider the function $f(x)=e^{-x}\cdot\sin x$, then
\begin{qoptions}
\opt{A}{Area bounded by the curve $f(x)$ and $x$-axis between ordinate $x=-2\pi$ and $x=0$ is $\left(\dfrac{e^{\pi}+1}{\sqrt2}\right)^2$}
\opt{B}{Area bounded by the curve $f(x)$ and $x$-axis between ordinate $x=-\pi$ and $x=\pi$ is $\dfrac{e^{\pi}-e^{-\pi}}{2}$}
\opt{C}{$\displaystyle\int_{m\pi}^{(m+1)\pi}f(x)\,dx<\int_{n\pi}^{(n+1)\pi}f(x)\,dx$ where $(m,n\in I)$ and $(m<n)$}
\opt{D}{$\displaystyle\int_{2m\pi}^{(2m+1)\pi}f(x)\,dx<\int_{2n\pi}^{(2n+1)\pi}f(x)\,dx$ where $(m,n\in I)$ and $(m>n)$}
\end{qoptions}

\qn{3} Consider differential equation\[\begin{aligned}&y^2\left(x\cdot e^y\,dy-e^y\cdot dx-x\cdot dx\right)\\&\quad+xy\left(x\,dy-\left(\dfrac{dy}{y}+\dfrac{dx}{x}\right)\right)=x^2y^2\end{aligned}\]

If $y(1)=1$. If solution of differential equation is $\dfrac{1}{x}(h(y))+\ln(g(x,\ y))=x+e$, then
\begin{qoptions}
\opt{A}{$h(2)=\dfrac{2\cdot e^2+1}{2}$}
\opt{B}{$g(x,y)=5$ represent a parabola passing through origin}
\opt{C}{$h(x)$ is an increasing function for $x\in[1,\infty)$}
\opt{D}{$g(x,y)=5$ represent a straight line}
\end{qoptions}

\qn{4} $\displaystyle\lim_{n\to\infty}\dfrac{\left(1^\lambda+2^\lambda+3^\lambda+\ldots+n^\lambda\right)}{\left(1^4+2^4+\ldots+n^4\right)\left(1^3+2^3+3^3+\ldots+n^3\right)}=F(\lambda),\ \lambda\in N$ then
\begin{qoptions}
\opt{A}{$F(\lambda)$ is finite for $\lambda\le8$}
\opt{B}{$F(8)=\dfrac{20}{11}$}
\opt{C}{$F(8)=\dfrac{20}{9}$}
\opt{D}{$F(7)=0$}
\end{qoptions}

\qn{5} If the least value of area bounded by $y=x^2-3$ and $y=kx+2$ be $A$, then the value of $\dfrac{3A}{10\sqrt5}$ is
\begin{qoptions}
\opt{A}{less than 2}
\opt{B}{greater than 5}
\opt{C}{less than 4}
\opt{D}{greater than 1}
\end{qoptions}

\qn{6} Let $f(x)$ be a polynomial of degree 3 such that $f(0)=1$, $f(1)=2$, $x=0$ is a critical point but $f(x)$ does not have local extremum at $x=0$, then $\displaystyle\int\dfrac{f(x)}{\sqrt{x^2+7}}dx$ is equal to
\begin{qoptions}
\opt{A}{$\dfrac{\sqrt{x^2+7}}{3}\left(x^2-14\right)+\ln\left(x+\sqrt{x^2+7}\right)+C$}
\opt{B}{$\dfrac{\sqrt{x^2+7}}{3}\left(x^2-14\right)-\ln\left(\sqrt{x^2+7}-x\right)+C$}
\opt{C}{$\ln\left(x+\sqrt{x^2+7}\right)+\sqrt{x^2+7}\left(x^2-7\right)+C$}
\opt{D}{$\ln\left(x-\sqrt{x^2+7}\right)+\sqrt{x^2+7}\left(x^2-7\right)+C$}
\end{qoptions}

\sect{SECTION-I(ii)}
\paraq{7}{8}
Let $f$ be a monotonic and differentiable function satisfying $\displaystyle\int_0^{f(x)}f^{-1}(t)\,dt-\int_0^{x}(\cos t-f(t))\,dt=0$ and $f(0)=1$.

\qn{7} The number of solution(s) of the equation $\begin{vmatrix}f(2x)&f(x)\\ \sin x&2\end{vmatrix}=0$ in $(0,2\pi)$ is
\begin{qoptions}
\opt{A}{3}
\opt{B}{6}
\opt{C}{4}
\opt{D}{5}
\end{qoptions}

\qn{8} If the value of $\displaystyle\int_0^{\pi/2}f(x)dx=\lambda$, then $[|\lambda|]$ is equal to, $[\ ]=$ G.I.F.
\begin{qoptions}
\opt{A}{2}
\opt{B}{3}
\opt{C}{4}
\opt{D}{1}
\end{qoptions}

\paraq{9}{10}
Let : $\displaystyle\int_0^{\pi}\dfrac{x}{1-\cos\phi\cdot\sin x}dx=f(\pi-\phi)$

\qn{9} Then the value of $f\left(\dfrac{\pi}{6}\right)$ is
\begin{qoptions}
\opt{A}{$\pi^2$}
\opt{B}{$\dfrac{\pi^2}{2}$}
\opt{C}{$\dfrac{\pi^2}{3}$}
\opt{D}{$\dfrac{\pi^2}{4}$}
\end{qoptions}

\qn{10} Then the value of $\displaystyle\int_0^{\pi/2}f^{\,2}(x)\,dx$ is
\begin{qoptions}
\opt{A}{$\pi^3\ln2$}
\opt{B}{$2\pi^3\ln2$}
\opt{C}{$\dfrac{\pi^3}{2}\ln2$}
\opt{D}{$3\pi^3\ln2$}
\end{qoptions}

\sect{SECTION-II(i)}
\paraq{1}{2}
Let $f(x)$ be a non constant function such that $f(x)=x^2-\displaystyle\int_0^1(f(t)+x)^2dt$
\begin{center}
$a=\displaystyle\int_0^\infty\dfrac{\tan^{-1}(x)}{x(1+x^2)}dx$\hfill $b=\displaystyle\int_0^1\dfrac{5x^8+3x^4+1}{\sqrt{x^8+x^4+1}}dx$\\[6pt]
$c=\displaystyle\int_0^2\sqrt{x^3+1}\,dx+\int_1^3\left(x^2-1\right)^{1/3}dx$\hfill $d=\displaystyle\int_0^1\dfrac{3x^2+9x+7}{(x^2+3x+2)^3}dx$
\end{center}

\qn{1} If $A$ is the area of region bounded by $f(x)$ and $x$ axis from $x=\dfrac{2a}{\pi\ln2}$ to $x=b^2$, then $\dfrac{A}{4}$ is

\qn{2} If $S$ area of region bounded by $f^{-1}(x)$ and $x$ axis from $x=\dfrac{72d}{11}$ to $x=c$, then $\dfrac{[S]}{2}$ is (where $[.]$ GIF)

\paraq{3}{4}
$I=\displaystyle\int f\left(x^m(ax^n+b)^{r/s}\right)dx;\ r,s\in I\ \&\ m,n\in Q.$

If $\dfrac{m+1}{n}=$ integer, then we can solve by substituting $ax^n+b=t^2$ and if $\dfrac{m+1}{n}\ne$ integer but $\dfrac{m+1}{n}+\dfrac{r}{s}=$ integer; we can substitute $a+bx^{-n}=t^s$; it converts function into integrable form.

\qn{3} If $\displaystyle\int\dfrac{dx}{\sqrt{a^{5x}}\sqrt[4]{\left(a^{2x}+a^{-2x}\right)^3}}=\dfrac{-1}{\ln a}g(x)+C$, then $g\left(-\dfrac14\log_a15\right)$ is equals -

\qn{4} $\displaystyle\int\sqrt{\operatorname{cosec}x}\left(2\sqrt[3]{\operatorname{cosec}x}+1\right)^{3/2}\cdot\cos x\,dx=pt^3+qt+r\sqrt2\ln\left|\dfrac{t-\sqrt2}{t+\sqrt2}\right|+C$ (where $t=\sqrt{2+\sin^{\frac13}x}$ )

then $p+q-r$ is equal

\paraq{5}{6}
Let $f(x)$ be a polynomial of degree 3 passing through origin and have a local maximum at $x=\dfrac{-1}{\sqrt{10}}$. Also $f'(x)$ has a local minimum at $x=0$ and $f(1)=7$.

\qn{5} The value of $f(3)$ is equal to

\qn{6} Consider a real valued function $y=g(x)$ satisfying the differential equation $\dfrac{dy}{dx}+f'(x)\cdot y=(f(x)+1)f'(x)$ such that $g(0)=1$ and $g(1)=\alpha+e^{-\beta}$. Then $(\alpha+\beta)$ is equal to

\sect{SECTION-II(ii)}
\qn{7} The ratio of the areas in which the curve $y=\left[\dfrac{x^3}{101}+\dfrac{x}{36}\right]$ divides the circle $x^2+y^2-4x+2y+1=0$, (where $[.]$ denotes greatest integer function) is $\dfrac{k\pi-3\sqrt3}{m\pi+3\sqrt3}$ then value of $|k+m|$ \underline{\hspace{2.5em}}

\qn{8} Consider the family of curve $y=f(x)$ whose differential equation is given by $x\dfrac{dy}{dx}+y=\dfrac{2\cos x+3}{(2+3\cos x)^2}$.

If $f\left(\dfrac{\pi}{2}\right)=\dfrac{2}{\pi}$, then value of $\left[\pi f\left(\dfrac{\pi}{3}\right)\right]$ is (where $[.]$ is GIF)

\qn{9} Let $f$ be a continuous function on $[0,\infty)$ such that $\displaystyle\lim_{x\to\infty}\left(f(x)+\int_0^xf(t)\,dt\right)$ exists ; then find $\displaystyle\lim_{x\to\infty}\left(f(x)+x\sin\dfrac1x\right)$
\end{multicols}
\end{document}
